% =====================================================================
%  Выводы по лабораторной работе.
% =====================================================================

\labpart{Выводы}

В ходе выполнения лабораторной работы разработана информационно-поисковая
система по варианту 22: собственный поисковый робот наполняет коллекцию
англоязычными документами, модуль индексирования строит поисковый образ
каждого документа по формулам инверсной частоты термина и веса термина
(1.5--1.6), а поиск по естественно-языковому запросу реализован через
векторную модель -- построение поискового образа запроса и ранжирование
документов по косинусной мере. Дополнительно к обязательной модели на
TF-IDF реализована вторая модель поиска -- на плотных векторных
представлениях, вычисляемых моделью Qwen3-Embedding-8B, -- для
последующего сравнения качества обоих подходов. Все ключевые формулы
вынесены в переиспользуемую библиотеку в виде чистых, покрытых тестами
функций, что позволило проверить их корректность независимо от
инфраструктуры (БД, HTTP).

Реализована и опробована на собранной коллекции \texttt{GTA} (1000
документов) оценка качества по методике ROMIP: Precision/Recall/F1 на
топ-$K$, Average Precision, R-Precision и интерполированная кривая
Precision/Recall. По 5 тестовым запросам дополнительная модель на
плотных векторных представлениях (Qwen3-Embedding-8B) превзошла
обязательную модель на TF-IDF по всем агрегированным метрикам (MAP
0{,}842 против 0{,}442, R-Precision 0{,}775 против 0{,}475) -- в
основном за счёт устойчивости к перефразировкам запроса, которые TF-IDF
не находит из-за отсутствия общих слов с текстом документа. При этом на
одном из пяти запросов, дословно совпадающем с формулировками в тексте
документов, точнее оказался именно TF-IDF -- то есть ни одна из моделей
не является безусловно лучшей, и разумным следующим шагом является
гибридное ранжирование, комбинирующее обе оценки. Также подтвердилось,
что более раннее исправление повторного использования одинаковых по
тексту запросов действительно снимает вырождение Recall в 1 на
достаточно большой коллекции. Набор тестовых запросов (5) при этом
всё ещё меньше рекомендованных методикой 15--20, что ограничивает
статистическую надёжность найденных различий.
